可微位姿估计层位于学习式前端与几何后端之间，必须对一个全局最优位姿求导，而该位姿在分支交叉处是不连续的。基于半定规划（SDP）松弛的确定性层在松弛紧时具有可认证性，但需要把问题提升为矩阵变量，并且在两个全局分支交叉时没有定义良好的行为。基于 Boltzmann 分布的概率层避免了硬性分支选择，但通常依赖六维位姿空间中的 Monte Carlo 积分。本文表明，二次位姿估计问题（QPEP）的四元数表述为第三种方案提供了代数结构。对于精确消去平移后剩余的旋转子问题，我们证明：标准 $10$ 维旋转提升恰好等于正交矩阵 $U$ 作用下的 $2Uv_2(q)$；21 条 $SO(3)$ 约束张成的空间恰为 Veronese 簇 $\nu_2(\mathbb{RP}^3)$ 的二次消没理想；相应的 $10\times10$ SDP 是 $\mathbb{S}^3$ 上四次型的 Veronese（Shor）松弛；在一个构造的 Choi--Lam 实例上，认证间隙为 $0.034188$，$\mathrm{rank}(X^\star)=6$。我们推导出闭式 Euler 乘子 $\lambda^\star=4J$，以及 $\mathbb{S}^3$ 上具有精确对映不变性的 $3\times3$ 黎曼隐式微分。将硬 argmin 替换为多模态 Laplace 后验——平移被精确积掉，旋转积分在所有 QPEP 极小点上解析近似，权重为 $w_k\propto e^{-J_k/\tau}/\sqrt{\det H_k}$——得到的层没有 Monte Carlo 方差，推理时无 SDP 提升，反向传播中每个活跃模态只需一次 $3\times3$ 求解。在真实的 LINEMOD 双模态实验中，硬 argmin 的中位旋转误差为 $2.06^\circ$、ADD-S 为 $3.7$ mm，而 Laplace 后验均值被拉向竞争模态（$44.24^\circ$、$13.9$ mm，成对偏差 $+36.9^\circ$）；全次数同伦在 $4/4$ 个实例上确认了通用的 $40$ 个复射影根。
